% Research-only exact identity; standard Gauss/CRT ingredients.
\subsection{Derivation and boundary of the joint phase cancellation}
\label{jdq:details}
Keep the notation of Section~\ref{jdq:core}. The actual factor count is
$r_U(m)$ and the common possibly complex profile is $W$. After fixing
$u_0$, let $F$ be the exact finite weight on $\mathcal O^6$ containing
all literal row restrictions. Only $v_1,\ldots,v_5$ are squarefree and
pairwise coprime; $v_6$ is arbitrary and retains its permitted overlaps.
No independence, factorization or smoothing of $F$ is assumed. A
restriction depending on the column orientation must instead be
partitioned explicitly or retained in an orientation-indexed transform;
the common-weight Parseval step below then requires separate checking.

CRT gives, for $1\le j\le5$ and $(k,q)=1$,
\[
 G_j(k)=\theta_{r,t}(\alpha)^j
       \tau_\beta(\theta_{r,t}^j)
       \overline{\theta_{r,t}(k)^j}\,c_h(k),
 \qquad G_6(k)=c_L(k).
\]
Indeed the $q$-frequency is $k\alpha^{-1}$ and the $h$-frequency is
$k\beta^{-1}$. The former contributes $\theta(\alpha)^j$;
the latter leaves $c_h(k)$ unchanged under unit scaling. A nonunit
$q$-frequency gives zero by primitivity. The usual exact Gauss identity
$\tau_\beta(\eta)\tau_\beta(\bar\eta)=\eta(-1)Q$ pairs powers
$1,5$ and $2,4$. Therefore
\[
 \prod_{j=1}^5\tau_\beta(\theta^j)
 =Q^2\theta(-1)^3\tau_\beta(\theta^3)
 =Q^2\xi_q(-1)\tau_\beta(\xi_q).
\]
Multiplying the five CRT factors contributes
$\theta(\alpha)^{15}=\xi_q(\alpha)$, proving
\eqref{jdq:gauss}. Its modulus is $Q^{5/2}$ times the displayed
Ramanujan magnitudes. Generator changes require corresponding changes
to the additive characters; the formulas are covariant under doing so.

Define the unnormalized finite Fourier transform
\[
 \widehat F_L(k)=\sum_{v\in\mathcal O^6}F(v)
             e_{\alpha\beta}\!\left(-\sum_{j=1}^6k_jv_j\right).
\]
Inversion of the periodization of $F$ has factor $\ell^{-6}$.
For $t=q/r$ put
\begin{align*}
 a_{h,q}(r;u_0)={}&r_U(hr)r_U(ht)\nu(r)\overline{\nu(t)}
 W(Nh\,Nr/D)\overline{W(Nh\,Nt/D)}\theta_{r,t}(u_0),\\
 A_{h,q}(K;u_0)={}&\sum_{r\mid q}a_{h,q}(r;u_0)
                         \overline{\theta_{r,q/r}(K)},\\
 B_{h,q,F}(K)={}&
 \sum_{\substack{k\bmod L,\ (k_1\cdots k_5,q)=1\\K(k)=K}}
 \widehat F_L(k)\prod_{j=1}^5c_h(k_j)c_L(k_6).
\end{align*}
Set $a_{h,q}=0$ for every excluded orientation, retaining the original
annuli, column support and actual long-long factor counts.
The common $\nu(h)$ phase cancels since $\nu$ is unitary on good ideals.
The normalized original C--C contribution of any chosen off-diagonal
sector $\mathcal T$ is exactly
\begin{equation}\label{jdq:joint}
 \frac1D\sum_{(h,q)\in\mathcal T}
 \mu(q)\frac{Q^2\xi_q(-\alpha)\tau_\beta(\xi_q)}{\ell^6}
 \sum_{K\in(\mathcal O/q)^\times}B_{h,q,F}(K)A_{h,q}(K;u_0).
\end{equation}
Sum all original $u_0$ classes and literal blocks afterwards.
This applies in particular to $h=1$ and to the original
low-gcd/large-conductor sector. No signed pieces have been discarded.
For C--$\mu$, replace the second $r_U(ht)$ by $1$ and change nothing else.

If $r\ne s$, the quotient of their orientation characters has
exponents $\pm2$ at the flipped primes, hence is a nontrivial cubic
character. Unit-group orthogonality proves
\eqref{jdq:orientation}, and so
\[
 \sum_K|A_{h,q}(K)|^2=\varphi(q)\sum_{r\mid q}|a_{h,q}(r)|^2.
\]
One normalized conductor term in \eqref{jdq:joint} consequently has the
finite bound
\begin{equation}\label{jdq:budget}
 |E_{h,q}|\le\frac{Q^{5/2}}{D\ell^6}
 \left(\varphi(q)\sum_{r\mid q}|a_{h,q}(r)|^2\right)^{1/2}
 \|B_{h,q,F}\|_2.
\end{equation}
The product map from five unit frequencies to $K$ has every fiber of
cardinality $\varphi(q)^4$, because $k_1$ is uniquely solvable.
All orientations agree pointwise exactly on
$\{K:x_p(K)^2=1\text{ for every }p\mid q\}$, of cardinality
$\varphi(q)/3^{\omega(q)}$; the common value is $x_q(K)\in\{\pm1\}$.
Thus neither uniform orthogonality nor five-frequency counting is an
independence statement for the true $\widehat F_L$ weight.

Centered fundamental-parallelogram representatives satisfy
$|k_j|\le|\alpha\beta|=\sqrt\ell$. This gives exact conductor-sized
finite completion. For any frequency subset $\Omega$, restrict $B$ to
$B_\Omega$. The error in one conductor term is at most the right side
of \eqref{jdq:budget} with $\|B-B_\Omega\|_2$ in place of $\|B\|_2$.
There is no proved small tail below that height: the smooth column
profile does not remove the original sharp row, ray or squarefree masks.

No $A_3$ assumption entered. In particular the full face
$A_1=A_2=A$, $AA_3=H^{1/3}$, $A_4=A_5=A_6=1$ and its neighborhood
are retained in $F$. On the pure $A_3$ endpoint where the other literal
factors equal $1$, inverse summation in their frequencies leaves
\[
 \left(\sum_{r\mid q}a_{h,q}(r)\right)
 \sum_vF_3(v)\xi_q(v)\mathbf1_{(v,h)=1}.
\]
There is no orientation cancellation in that row character.
On the general face keep $\theta(v_1)\theta(v_2)^2\xi_q(v_3)$
and all coupled masks. If only powers $1,2,3$ are completed after the
others are fixed, their Gauss product is
\[
 \tau(\theta)\tau(\theta^2)\tau(\theta^3)
 =\xi_q(-1)QJ(\theta,\theta^2),\qquad |J(\theta,\theta^2)|=\sqrt Q,
\]
whose phase can depend on orientation. Complete five-power phase
balance therefore does not remove every actual near-face phase.
The remaining task is a quantitative signed bound for the actual
$h,q,K$ sum, especially $h=1$, with $Q\asymp D^2/(Nh)^2$.
No estimate of $B$ sufficient for a new global exponent is supplied.
The old small-conductor/high-gcd sector is unchanged.

An exact standard-library checker under
\path{planning/joint-dual-20261105/research/} passes $991$ assertions:
local split/inert Gauss products, all $546$ complete raw transforms for
$q=(3+\omega)$, $h=(4+\omega)$, $L=(11+6\omega)$, twelve mixed
frequency products, the four orientations at conductor norm $91$,
orientation Parseval, the coherent cubic diagonal and $7776$ product
fibers. The receipt is
\path{outputs/checks/2026-11-05-joint-dual/research/joint-gauss-check-v1.json}.
These finite diagnostics do not prove an incomplete-moment exponent,
$BT_\sigma$, a new Lean theorem, zero location or GRH.
